\documentclass[a4paper]{article}

\usepackage{kdiss_conference_autumn}% autumn conference proceeding style file
\usepackage{kotex}
\usepackage{amsmath,graphicx,chicago}
\usepackage{amssymb}
\usepackage{multirow}
\usepackage{array}
\usepackage{hhline}

%% ==================== 여기는 고치지 마십시오. =======================
\submit{final}
\thefirstpage{1}
\thelastpage{1}
%% ====================================================================

\begin{document}

\title{비정상 시계열 변동성 예측: 순차 대 병렬 알고리즘 성능 비교}

\hauthor{윤태준$^{*}$~$\cdot$~김채원$^{*}$~$\cdot$~손낙훈$^{*}$}

\haddress{${}^{*}$계명대학교 자연과학대학 통계학과}

\begin{abstract}
암호화폐 가격은 평균 수준이 일정하지 않고 변동성이 한동안 높거나 낮게 유지되다가 갑자기 바뀌는 비정상 시계열이다.
이런 자료에서 평균 성능 하나로 예측 모델을 고르면 변동이 클 때 적합하지 않은 모델을 고를 수 있다.
본 연구는 업비트 20개 종목의 15분 단위 가격(2023년 10월$\sim$2026년 10월)으로 15분부터 12시간 뒤까지의 변동성을 예측하고,
과거 시점을 차례로 읽는 순차 모델(GARCH, 순환 신경망)과 과거 구간 전체를 한꺼번에 처리하는 병렬 모델(트리 기반 모델, 커널 회귀, 트랜스포머, 사전학습 시계열 모델) 27개의 성능을 비교하였다.
예측 시점의 직전 변동성 크기로 자료를 다섯 단계로 나누고, 모델 간 예측 오차의 차이가 통계적으로 유의한지 검정하였다.
그 결과 변동이 작거나 보통일 때는 트리 기반 모델과 커널 회귀가, 변동이 가장 클 때 1시간 이상 예측에서는 순환 신경망이 가장 작은 오차를 보였다.
다만 변동이 가장 클 때 순환 신경망과 트리$\cdot$커널 모델의 차이는 통계적으로 유의하지 않았다.
같은 단계에서 병렬 딥러닝 모델은 모든 비교에서 최선 모델보다 유의하게 나빴고, 순환 신경망은 병렬 딥러닝 모델과 같은 입력을 주어도 더 나은 성능을 유지하였다.
이 결과는 비정상 시계열의 변동성 예측에서 하나의 최적 모델을 찾기보다 변동성 상황에 따라 모델을 선택해야 함을 보여 준다.
\end{abstract}

\keywords{비정상 시계열, 변동성 예측, 순환 신경망, 트랜스포머, 암호화폐.}

\end{document}
